\subsection{COMPACT GROUPS OF RANK 1}

By General Topology, Chap. VIII, §1, no. 4, Prop. 3, Prop. 4 and Remark 4, the topological group \(\mathbf{SU}(2,\mathbf{C})\) is isomorphic to the topological group \(\mathbf{S}_3\) of quaternions of norm 1, and the quotient of \(\mathbf{SU}(2,\mathbf{C})\) by the subgroup Z consisting of the matrices \(I_2\) and \(-I_2\) is isomorphic to the topological group \(\mathbf{SO}(3,\mathbf{R})\) . Note that Z is the centre of \(\mathbf{SU}(2,\mathbf{C})\) : indeed, since \(\mathbf{H} = \mathbf{R}.\mathbf{S}_3\) , every element of the centre of the group \(\mathbf{S}_3\) is in the centre \(\mathbf{R}\) of the algebra \(\mathbf{H}\) and hence belongs to the group with two elements \(\mathbf{S}_3 \cap \mathbf{R} = \{-1,1\}\) .

PROPOSITION 6. Every compact semi-simple real Lie algebra of rank 1 is isomorphic to \(\mathfrak{su}(2,\mathbf{C})\) and to \(\mathfrak{o}(3,\mathbf{R})\) . Every connected semi-simple compact Lie group of rank 1 is isomorphic to \(\mathbf{SU}(2,\mathbf{C})\) if it is simply-connected, and to \(\mathbf{SO}(3,\mathbf{R})\) if not.

The first assertion follows from the Cor. of Prop. 4 and Prop. 5. Since \(\mathbf{SU}(2,\mathbf{C})\) is homeomorphic to \(\mathbf{S}_3\) (General Topology, Chap. VIII, §1, no. 4, Remark 4), hence is simply-connected (General Topology, Chap. XI, in preparation), every simply-connected compact semi-simple Lie group of rank 1 is isomorphic to \(\mathbf{SU}(2,\mathbf{C})\) ; every connected compact semi-simple Lie group of rank 1 that is not simply-connected is isomorphic to a quotient of \(\mathbf{SU}(2,\mathbf{C})\) by a subgroup of \(Z\) that does not reduce to the identity element, hence to \(\mathbf{SO}(3,\mathbf{R})\) .

Remark. We have seen above that \(\mathbf{SU}(2,\mathbf{C})\) is simply-connected and that \(\pi_{1}(\mathbf{SO}(3,\mathbf{R}))\) is of order 2. We shall see much later that these results generalize to \(\mathbf{SU}(n,\mathbf{C})\) ( \(n\geq1\) ) and \(\mathbf{SO}(n,\mathbf{R})\) ( \(n\geq3\) ), respectively (cf.~also §3, Exerc. 4 and 5).

Recall (Chap. VIII, §1, no. 1) that the canonical basis of \(\mathfrak{sl}(2,\mathbf{C})\) is the basis \((X_{+},X_{-},H)\) , where

\[
X _ {+} = \left( \begin{array}{c c} 0 & 1 \\ 0 & 0 \end{array} \right), X _ {-} = \left( \begin{array}{c c} 0 & 0 \\ - 1 & 0 \end{array} \right), H = \left( \begin{array}{c c} 1 & 0 \\ 0 & - 1 \end{array} \right).
\]

We thus obtain a basis \((U, V, iH)\) of \(\mathfrak{su}(2, \mathbf{C})\) , also called canonical, by putting

\[
\begin{array}{c} U = X _ {+} + X _ {-} = \left( \begin{array}{c c} 0 & 1 \\ - 1 & 0 \end{array} \right), V = i (X _ {+} - X _ {-}) = \left( \begin{array}{c c} 0 & i \\ i & 0 \end{array} \right), \\ i H = \left( \begin{array}{c c} i & 0 \\ 0 & - i \end{array} \right). \end{array}
\]

We have

\[
[ i H, U ] = 2 V, [ i H, V ] = - 2 U, [ U, V ] = 2 i H.\tag{13}
\]

If B denotes the Killing form of \(\mathfrak{su}(2,\mathbf{C})\) an immediate calculation gives

\[
\mathrm{B} (a U + b V + c i H, a ^ {\prime} U + b ^ {\prime} V + c ^ {\prime} i H) = - 8 (a a ^ {\prime} + b b ^ {\prime} + c c ^ {\prime}),\tag{14}
\]

so that, if we identify \(\mathfrak{su}(2,\mathbf{C})\) with \(R^{3}\) by means of the canonical basis, the adjoint representation of \(\mathbf{SU}(2,\mathbf{C})\) defines a homomorphism \(\mathbf{SU}(2,\mathbf{C})\to\mathbf{SO}(3,\mathbf{R})\) (cf.~above).

Further, note that \(\mathbf{Ri}H\) is a Cartan subalgebra of \(\mathfrak{su}(2,\mathbf{C})\) , that the maximal torus \(T\) of \(\mathbf{SU}(2,\mathbf{C})\) that corresponds to it consists of the diagonal matrices \(\left( \begin{array}{cc}a & 0\\ 0 & \bar{a} \end{array} \right)\) , where \(a\bar{a} = 1\) , and that the exponential map

\(\exp : \mathbf{R}iH \to \mathrm{T}\)

maps \(xH\) , for \(x \in \mathbf{R}i\) , to the matrix \(\left( \begin{array}{cc}\exp (x) & 0\\ 0 & \exp (-x) \end{array} \right)\) , and thus has kernel \(\mathbf{Z}.K\) , where \(K\) is the element of \(\mathfrak{su}(2,\mathbf{C})\) defined by

\[
K = 2 \pi i H = \left( \begin{array}{c c} 2 \pi i & 0 \\ 0 & - 2 \pi i \end{array} \right).\tag{15}
\]

Further, the centre of \(\mathbf{SU}(2,\mathbf{C})\) consists of the identity and \(\exp (K / 2)\) .

\[
\theta = \left( \begin{array}{c c} 0 & - 1 \\ 1 & 0 \end{array} \right) \in \mathbf {S U} (2, \mathbf {C}).\tag{16}
\]

By Chap. VIII, §1, no. 5,

\[
\theta^ {2} = \left( \begin{array}{c c} - 1 & 0 \\ 0 & - 1 \end{array} \right), (\text {Int}   \theta) t = t ^ {- 1}, t \in \text {T},\tag{17}
\]

\[
(\operatorname{Ad} \theta) X _ {+} = X _ {-}, (\operatorname{Ad} \theta) X _ {-} = X _ {+}, (\operatorname{Ad} \theta) U = U, (\operatorname{Ad} \theta) V = - V.\tag{18}
\]

\[
\text {   Finally,   for   } t = \left( \begin{array}{c c} a & 0 \\ 0 & \bar {a} \end{array} \right) \in \mathrm{T}, \text {   we   have   }
\]

\[
(\operatorname{Ad} t) X _ {+} = a ^ {2} X _ {+}, (\operatorname{Ad} t) X _ {-} = a ^ {- 2} X _ {-}, (\operatorname{Ad} t) H = H,\tag{19}
\]

\[
(\operatorname{Ad} t) U = \mathcal {R} (a ^ {2}) U + \mathcal {I} (a ^ {2}) V, (\operatorname{Ad} t) V = - \mathcal {I} (a ^ {2}) U + \mathcal {R} (a ^ {2}) V.\tag{20}
\]

